% Uebungsblatt Lineare Algebra (German, university).
% Headings: "Aufgabe 6.1 (Titel, P Punkte)" in bold, sheet.task numbering; parts a), b), ...
\documentclass[11pt,a4paper]{article}
\usepackage[T1]{fontenc}
\usepackage[utf8]{inputenc}
\usepackage{lmodern}
\usepackage[ngerman]{babel}
\usepackage[margin=2.3cm,top=2.0cm,bottom=2.2cm]{geometry}
\usepackage{amsmath,amssymb}
\usepackage{enumitem}

\setlength{\parindent}{0pt}
\setlist[enumerate,1]{label=\alph*), leftmargin=2em, itemsep=0.35em, topsep=0.3em}
\newcommand{\aufgabe}[3]{\bigskip\par\textbf{Aufgabe #1 (#2, #3 Punkte)}\par\nopagebreak\smallskip}
\newcommand{\R}{\mathbb{R}}

\begin{document}

% ---- Kopf -----------------------------------------------------------------------------
{\small Lineare Algebra I \textperiodcentered{} Wintersemester 2026/27 \hfill Mathematisches Institut}\\[-0.4em]
\rule{\linewidth}{0.6pt}\\[0.5em]
{\Large\bfseries Übungsblatt 6}\\[0.4em]
Abgabe: Montag, 16.\ November 2026, vor der Vorlesung \hfill Gesamtpunktzahl: 20

% ---- Aufgabe 6.1 ----------------------------------------------------------------------
\aufgabe{6.1}{Gauß-Verfahren}{6}
Gegeben ist das lineare Gleichungssystem über $\R$
\[
\begin{aligned}
  x + \phantom{2}y + \phantom{3}z &= \phantom{-}4,\\
  2x - \phantom{2}y + 3z &= 14,\\
  x + 2y - \phantom{3}z &= -3.
\end{aligned}
\]
\begin{enumerate}
  \item Bringen Sie die erweiterte Koeffizientenmatrix mit dem Gauß-Verfahren auf Zeilenstufenform.
  \item Bestimmen Sie die Lösungsmenge des Gleichungssystems.
  \item Berechnen Sie die Determinante der Koeffizientenmatrix. Was sagt sie über die Lösbarkeit des
        Systems aus?
\end{enumerate}

% ---- Aufgabe 6.2 ----------------------------------------------------------------------
\aufgabe{6.2}{Determinante}{4}
Sei
\[
A = \begin{pmatrix} 2 & 0 & 1 \\ 1 & 3 & 2 \\ 0 & -1 & 4 \end{pmatrix}.
\]
\begin{enumerate}
  \item Berechnen Sie $\det A$ durch Entwicklung nach der ersten Zeile. Ist $A$ invertierbar?
  \item Bestimmen Sie $\det(2A)$, $\det(A^{\top}A)$ und $\det(A^{-1})$, ohne die Matrizen
        auszurechnen.
\end{enumerate}

% ---- Aufgabe 6.3 ----------------------------------------------------------------------
\aufgabe{6.3}{Eigenwerte und Eigenvektoren}{6}
Gegeben sei die Matrix
\[
B = \begin{pmatrix} 4 & 1 \\ 2 & 3 \end{pmatrix} \in \R^{2 \times 2}.
\]
\begin{enumerate}
  \item Bestimmen Sie das charakteristische Polynom von $B$ und alle Eigenwerte.
  \item Bestimmen Sie zu jedem Eigenwert einen Eigenvektor.
  \item Ist $B$ diagonalisierbar? Geben Sie gegebenenfalls eine invertierbare Matrix $P$ und eine
        Diagonalmatrix $D$ mit $B = PDP^{-1}$ an.
\end{enumerate}

% ---- Aufgabe 6.4 ----------------------------------------------------------------------
\aufgabe{6.4}{Beweisaufgabe}{4}
Sei $A \in \R^{n \times n}$.
\begin{enumerate}
  \item Zeigen Sie, dass jeder Eigenwert von $A$ gleich $0$ oder $1$ ist, wenn $A^2 = A$ gilt.
  \item Zeigen Sie, dass $\lambda^{-1}$ ein Eigenwert von $A^{-1}$ ist, wenn $A$ invertierbar ist und
        $\lambda$ ein Eigenwert von $A$ ist.
\end{enumerate}

\end{document}
